%===================================================================================================
\section{Reporting: Formatting and Integration with the Test Pipeline} %++++++++++++++++++++++++++++
  \label{Reporting: Formatting and Integration with the Test Pipeline} %++++++++++++++++++++++++++++
%===================================================================================================


\improve


\begin{justify}
Assertify/Checkify are designed to integrate into existing test workflows by wrapping standard
frameworks and replacing baseline failure text with a structured report. The integration strategy
follows the ``minimal intrusion'' principle: tests are still executed by conventional runners, but
failures are intercepted and re-raised with enriched information in a stable format.
\end{justify}


%===============================================================================
\subsection{Integration with example-based unit tests} %++++++++++++++++++++++++
     \label{Integration with example-based unit tests} %++++++++++++++++++++++++
%===============================================================================


\begin{justify}
For example-based tests, Assertify acts as a wrapper around standard assertion-style checks by
evaluating the quoted test expression and failing the test if it evaluates to \texttt{false}.
When a failure occurs, Assertify produces a report that explicitly includes:

\begin{itemize}
  \item the tested/failed expression,
  \item an optional simplified expression rendering, and
  \item the expected and actual values.
\end{itemize}

This structure is illustrated by the existing before/after comparison listings in the draft, where
Assertify reports the fully qualified expression alongside a simplified rendering and then prints
expected and actual results in a dedicated block.
\end{justify}


%===============================================================================
\subsection{Integration with property-based tests} %++++++++++++++++++++++++++++
     \label{Integration with property-based tests} %++++++++++++++++++++++++++++
%===============================================================================


\begin{justify}
For property-based tests, Checkify provides a wrapper that executes the property via FsCheck and,
on failure, captures both the original and shrunk counterexample inputs. These inputs are
incorporated into the structured report to make the falsifying case reproducible and to ensure that
the student-facing expression is presented with the relevant (typically shrunk) input instantiated.
This preserves the value of shrinking while presenting the result in the same
expected/actual-oriented format as example-based tests.
\end{justify}


%===============================================================================
\subsection{Formatting strategy} %++++++++++++++++++++++++++++++++++++++++++++++
     \label{Formatting strategy} %++++++++++++++++++++++++++++++++++++++++++++++
%===============================================================================


\begin{justify}
The reporting format is designed around a high-signal summary that foregrounds the novice-relevant
fields. The draft output examples show explicit section markers (e.g., for ``TESTED EXPRESSION'' and
``EXPECTED AND ACTUAL''), which support scanning and reduce the likelihood that students miss the
crucial mismatch due to stack trace noise. Additional details such as stack traces remain available
for traceability but are treated as secondary information.
\end{justify}


%===============================================================================
\subsection{Configurable verbosity and optional context} %++++++++++++++++++++++
     \label{Configurable verbosity and optional context} %++++++++++++++++++++++
%===============================================================================


\begin{justify}
Finally, reporting is configurable with respect to optional context components. In line with the
stated design considerations, the library supports toggling whether reductions and history
information are included, allowing the same testing setup to be used with different levels of
verbosity depending on the educational context (introductory debugging vs.\ more advanced
troubleshooting).
\end{justify}


%===================================================================================================
% EOF ++++++++++++++++++++++++++++++++++++++++++++++++++++++++++++++++++++++++++++++++++++++++++++++
%===================================================================================================